% ============================================================================
% CAPÍTULO 1: EL DESCUBRIMIENTO
% ============================================================================
\chapter{El Descubrimiento}
\label{cap:descubrimiento}

\begin{center}
\textit{``Todo comenzó con el número 22''}
\end{center}

\vspace{0.5cm}

\begin{tcolorbox}[colback=white, colframe=kleinblue, boxrule=2pt]
\centering
\textit{``La ciencia avanza funeral a funeral.''}\\
--- Max Planck\\[0.3cm]
\textit{``...o cuando alguien nota un patrón que siempre estuvo ahí.''}\\
--- Teoría Klein
\end{tcolorbox}


% ============================================================================
\section{Un Número Misterioso: 22}
% ============================================================================

\begin{analogiabox}
\textbf{El Detective Numérico}

Imagina que eres un detective investigando una escena del crimen.
En cada habitación encuentras el número 22: en la cerradura, en el
reloj, en la temperatura. ¿Coincidencia? Quizás la primera vez.
¿Pero diez veces? ¿Veinte? Eso es una pista.

Así comenzó nuestra investigación. El número 22 aparecía
demasiadas veces en la física de ondas gravitacionales.
\end{analogiabox}

En 2015, LIGO detectó las primeras ondas gravitacionales de la historia.
Estas ondulaciones en el espacio-tiempo provenían de dos agujeros negros
fusionándose a miles de millones de años luz de distancia.

Al analizar los datos, notamos algo peculiar:

\begin{itemize}
    \item La frecuencia pico del evento GW150914: $\sim$22 Hz
    \item Relaciones de masa que involucraban el número 22
    \item Patrones temporales con múltiplos de 22
\end{itemize}

¿Por qué 22? No hay ninguna razón física obvia.

22 no es un número fundamental conocido. No es $\pi$, no es $e$, no es
la constante de estructura fina. Es simplemente... 22.

O eso pensábamos.


% ============================================================================
\section{La Coincidencia que Cambió Todo: $22 \approx 7\pi$}
% ============================================================================

\begin{nivelmatematico}
\textbf{La Primera Ecuación}

\begin{equation}
\boxed{22 \approx 7\pi}
\end{equation}

\begin{align}
7 \times \pi &= 7 \times 3.14159... = 21.991149...\\[0.3cm]
\text{Error} &= \frac{|22 - 7\pi|}{22} = 0.0403\%
\end{align}

¡Menos del 0.05\% de error!
\end{nivelmatematico}

¿Por qué esto es importante?

El número $\pi$ aparece en TODA la física: círculos, esferas, ondas, rotaciones.
El número 7... ese era el misterio.

Entonces recordé: la botella de Klein.

En topología, la botella de Klein es una superficie cerrada no orientable.
En ciertas parametrizaciones, tiene una estructura de 7 capas o regiones.

¿Y si 22 Hz no es arbitrario?\\
¿Y si es $7\pi$ Hz --- la frecuencia natural de un universo Klein?

\begin{nivelconceptual}
\textbf{La Hipótesis Klein}

\textbf{Hipótesis:} El espacio-tiempo tiene una estructura topológica
similar a una botella de Klein, con 7 ``capas'' o regiones.

Esta estructura determina las constantes fundamentales de la física
a través del factor de supresión:

\begin{equation}
\boxed{7\pi \approx 22}
\end{equation}
\end{nivelconceptual}


% ============================================================================
\section{El Método Científico Inverso}
% ============================================================================

Normalmente, la ciencia funciona así:

\begin{center}
Teoría $\to$ Predicción $\to$ Experimento $\to$ Verificación
\end{center}

Nosotros lo hicimos al revés:

\begin{center}
Observación (22) $\to$ Patrón ($7\pi$) $\to$ Teoría (Klein) $\to$ MÁS predicciones
\end{center}

Esto es válido. Así se descubrieron:

\begin{itemize}
    \item La tabla periódica (patrones en elementos)
    \item La radiación cósmica de fondo (ruido en antenas)
    \item Los quarks (patrones en partículas)
\end{itemize}

El test de una teoría no es cómo se origina, sino si hace predicciones
correctas sobre cosas NUEVAS.

Así que nos preguntamos:

\begin{center}
\textit{``Si $22 = 7\pi$ describe las ondas gravitacionales,\\
¿qué más puede predecir?''}
\end{center}


% ============================================================================
\section{Las Primeras Predicciones}
% ============================================================================

Comenzamos a buscar el factor $7\pi$ en otras constantes físicas.

Lo que encontramos fue asombroso:

\begin{nivelcodigo}
\textbf{Verificación de Predicciones Iniciales}

\textbf{1. Velocidad de la luz:}
\begin{align}
c &= \left(3 - \frac{1}{(7\pi)^2}\right) \times 10^8 \text{ m/s}\\
\text{Predicción:} & \quad 299,792,459 \text{ m/s}\\
\text{Observado:} & \quad 299,792,458 \text{ m/s}\\
\text{Error:} & \quad 0.0003\%
\end{align}

\textbf{2. Constante de estructura fina:}
\begin{align}
\frac{1}{\alpha} &= 7^2\pi - 7 - \pi^2\\
\text{Predicción:} & \quad 137.0683\\
\text{Observado:} & \quad 137.0360\\
\text{Error:} & \quad 0.024\%
\end{align}

\textbf{3. Ratio masa protón/electrón:}
\begin{align}
\frac{m_p}{m_e} &= 6\pi^5 = (7-1)\pi^5\\
\text{Predicción:} & \quad 1836.1181\\
\text{Observado:} & \quad 1836.1527\\
\text{Error:} & \quad 0.002\%
\end{align}
\end{nivelcodigo}

\vspace{0.5cm}

¡TRES predicciones con errores menores al 0.03\%!

Esto no puede ser coincidencia. La probabilidad de que tres números
aleatorios coincidan con la realidad con esta precisión es:

\begin{equation}
P \approx (0.0003) \times (0.0003) \times (0.00002) \approx 10^{-12}
\end{equation}

Una en un billón.

\begin{analogiabox}
\textbf{El Código Postal}

Imagina que alguien te da un código: ``7-3.14159''

Con ese código, puedes calcular:
\begin{itemize}
    \item La velocidad de la luz
    \item La fuerza electromagnética
    \item La masa del protón
\end{itemize}

¿No sería como si el universo tuviera un ``código postal''
que define todas sus propiedades?

Ese código es: $7\pi$
\end{analogiabox}


% ============================================================================
\section{Resumen del Capítulo}
% ============================================================================

\begin{tcolorbox}[colback=kleingray, colframe=kleinblue, title=\textbf{Resumen del Capítulo 1}]

\textbf{IDEA CENTRAL:}

El número 22, observado en ondas gravitacionales, es aproximadamente
igual a $7\pi$. Esto sugiere que el universo tiene una estructura
topológica de 7 capas (botella de Klein).

\vspace{0.3cm}
\textbf{ECUACIÓN FUNDAMENTAL:}

\begin{equation}
\boxed{22 \approx 7\pi}
\end{equation}

\vspace{0.3cm}
\textbf{PREDICCIONES VERIFICADAS:}

\begin{enumerate}
    \item $c = (3 - 1/(7\pi)^2) \times 10^8$ m/s \hfill [0.0003\% error]
    \item $1/\alpha = 7^2\pi - 7 - \pi^2$ \hfill [0.024\% error]
    \item $m_p/m_e = 6\pi^5$ \hfill [0.002\% error]
\end{enumerate}

\vspace{0.3cm}
\textbf{PRÓXIMO CAPÍTULO:}

¿Qué es exactamente una botella de Klein?\\
¿Por qué tiene 7 capas?\\
¿Cómo se conecta con la física?

\end{tcolorbox}


% ============================================================================
\section*{Ejercicios}
% ============================================================================

\begin{enumerate}
    \item Calcula $7\pi$ con tu calculadora. ¿Cuánto difiere de 22?

    \item Si el universo fuera una botella de Klein con 8 capas en lugar de 7,
    ¿cuánto sería $8\pi$? ¿Hay algún fenómeno físico asociado a $\sim$25?

    \item El factor 6 en $m_p/m_e = 6\pi^5$ es igual a $7-1$.
    ¿Por qué podría tener sentido restar 1 de las 7 capas?

    \item Investiga: ¿En qué otros contextos aparece el número 22 en física?
\end{enumerate}


% ============================================================================
\section*{Código de Verificación}
% ============================================================================

\begin{lstlisting}[caption={Verificación numérica - Capítulo 1}]
import numpy as np
from numpy import pi

# Constante Klein
siete_pi = 7 * pi
print(f"7*pi = {siete_pi:.10f}")
print(f"22   = 22.0")
print(f"Error = {abs(siete_pi - 22)/22*100:.6f}%")

# Velocidad de la luz
c_obs = 299792458
c_pred = (3 - 1/siete_pi**2) * 1e8
print(f"\nc observado  = {c_obs} m/s")
print(f"c prediccion = {c_pred:.0f} m/s")
print(f"Error = {abs(c_pred - c_obs)/c_obs*100:.6f}%")

# Estructura fina
alpha_inv_obs = 137.035999
alpha_inv_pred = 7**2 * pi - 7 - pi**2
print(f"\n1/alpha observado  = {alpha_inv_obs:.6f}")
print(f"1/alpha prediccion = {alpha_inv_pred:.6f}")
print(f"Error = {abs(alpha_inv_pred - alpha_inv_obs)/alpha_inv_obs*100:.6f}%")

# Ratio de masas
mp_me_obs = 1836.15267
mp_me_pred = 6 * pi**5
print(f"\nm_p/m_e observado  = {mp_me_obs:.5f}")
print(f"m_p/m_e prediccion = {mp_me_pred:.5f}")
print(f"Error = {abs(mp_me_pred - mp_me_obs)/mp_me_obs*100:.6f}%")
\end{lstlisting}
